%%%%%%%%%%%%%%%%%%%%%%%%%%%%%%%%%%%%%%%%%%%%%%%%%%%%%%%%%%%%%%%%%%%%
%% I, the copyright holder of this work, release this work into the
%% public domain. This applies worldwide. In some countries this may
%% not be legally possible; if so: I grant anyone the right to use
%% this work for any purpose, without any conditions, unless such
%% conditions are required by law.
%%%%%%%%%%%%%%%%%%%%%%%%%%%%%%%%%%%%%%%%%%%%%%%%%%%%%%%%%%%%%%%%%%%%

\documentclass{beamer}
\usetheme[faculty=ped,basePath=../fibeamer,logo=../fibeamer/logo/mu/Logo-UC-3.png]{fibeamer}
\graphicspath{{../resources/}}

\usepackage{algorithm}
\usepackage{algpseudocode}

\usepackage[utf8]{inputenc}
\usepackage[
  main=english, %% By using `czech` or `slovak` as the main locale
                %% instead of `english`, you can typeset the
                %% presentation in either Czech or Slovak,
                %% respectively.
  czech, slovak %% The additional keys allow foreign texts to be
]{babel}        %% typeset as follows:
%%
%%   \begin{otherlanguage}{czech}   ... \end{otherlanguage}
%%   \begin{otherlanguage}{slovak}  ... \end{otherlanguage}
%%
%% These macros specify information about the presentation
\title{Variables, operadores y comparadores} %% that will be typeset on the
\subtitle{Semestre II-2026} %% title page.
\author{Andreina Cota Vidaurre}
%% These additional packages are used within the document:
\usepackage{ragged2e}  % `\justifying` text
\usepackage{booktabs}  % Tables
\usepackage{tabularx}
\usepackage{tikz}      % Diagrams
\usetikzlibrary{calc, shapes, backgrounds}
\usetikzlibrary{positioning,arrows.meta, shadows.blur}
\usepackage{amsmath, amssymb}
\usepackage{url}       % `\url`s
\usepackage{listings}  % Code listings
\usepackage{graphicx}
\usepackage{amsmath}
\usepackage[T1]{fontenc}
\usepackage{xcolor}
\usepackage{minted}
% \usepackage{adjustbox}


% Definicion de pseudocodigo en espaniol
\lstdefinelanguage{PseudocodigoES}{
    morekeywords={INICIO,FIN,PARA,HACER,SI,ENTONCES,FIN_SI,FIN_PARA,IMPRIMIR},
    sensitive=true,
    morecomment=[l]{//},
    morestring=[b]"
}

% Definicion de pseudocodigo en ingles
\lstdefinelanguage{PseudocodeEN}{
    morekeywords={BEGIN,END,FOR,DO,IF,THEN,END_IF,END_FOR,PRINT},
    sensitive=true,
    morecomment=[l]{//},
    morestring=[b]"
}

\lstset{
    basicstyle=\ttfamily\small,
    keywordstyle=\bfseries,
    columns=fullflexible,
    keepspaces=true,
    showstringspaces=false,
    frame=single,
    xleftmargin=0em,
    framexleftmargin=0em,
    framesep=2pt,
    gobble=4,
    resetmargins=true,
    aboveskip=0pt,
    belowskip=0pt
}

\lstdefinestyle{pythonstyle}{
    language=Python,
    basicstyle=\ttfamily\small,
    keywordstyle=\bfseries,
    commentstyle=\itshape,
    stringstyle=\ttfamily,
    showstringspaces=false,
    columns=fullflexible,
    keepspaces=true,
    frame=single,
    breaklines=true,
    xleftmargin=0em,
    framexleftmargin=0em,
    framesep=2pt,
    aboveskip=0pt,
    belowskip=0pt,
    gobble=4,
    resetmargins=true,
    emph={print,input,int,float,str,bool,len,range},
    emphstyle=\bfseries
}

\lstset{style=pythonstyle}

% \newcolumntype{Y}{>{\raggedright\arraybackslash}X}

\frenchspacing
\begin{document}
  \shorthandoff{-}
  \frame[c]{\maketitle}

% \begin{darkframes}

    \begin{frame}[fragile,t]{Contenido}
        \begin{enumerate}
            \item Variables (int, float, string, bool)
            \item Operadores (\textcolor{cyan}{$+, -, *, /, **, //, \%$})
            \item Comparadores (\textcolor{cyan}{$<, >, <=, >=, ==, !=$})
        \end{enumerate}
    \end{frame}

    \begin{frame}{Variables, operadores y comparadores}
         \begin{itemize}
             \item Un programa necesita tres cosas básicas:
             \begin{enumerate}
                 \item \textbf{Variables} para guardar datos
                 \item \textbf{Operadores} para trabajar con esos datos
                 \item \textbf{Comparadores} para revisar condiciones y decidir qué hacer
             \end{enumerate}
         \end{itemize}

         ``Primero el programa recuerda, luego calcula y finalmente decide``
    \end{frame}

    \begin{frame}{Variables, operadores y comparadores}
         \begin{itemize}
             \item Guardamos datos: edad, nombre, saldo, temperatura
             \item Hacemos operaciones: sumar compras, calcular promedio, restar gastos
             \item Comparamos: si una nota es mayor a 4, si hay suficiente saldo, si una clave coincide
         \end{itemize}
    \end{frame}

    \begin{frame}{Datos e información}
        \begin{itemize}
            \item Un \textbf{dato} es un valor aislado, sin mucho contexto por si solo
            \item La \textbf{información} aparece cuando los datos tienen \textbf{contexto} y adquieren significado para una persona 
        \end{itemize}
        Ejemplos:
        \begin{itemize}
            \item $25$ -> dato
            \item $edad = 25$ -> Tiene más significado
            \item ``La edad del estudiante es 25 años`` -> información
        \end{itemize}
    \end{frame}

    \begin{frame}[fragile,t]{Variables}
        \begin{itemize}
            \item Una variable almacena datos
            \item Tiene un nombre
            \item Guarda un valor
            \item Ese valor puede cambiar durante el programa
        \end{itemize}
        \centering
        \begin{minted}{python}
            edad = 18
            nombre = "Ana"
            promedio = 76.5
        \end{minted}
    \end{frame}

    \begin{frame}[fragile,t]{Variables}
        En matemáticas:
        \begin{itemize}
            \item $x$ representa un valor desconocido o general
            \item Ejemplo: $x + 3 = 10$
        \end{itemize}
        En programación:
        \begin{itemize}
            \item una variable guarda un valor concreto en memoria.
            
            Ejemplo:
            
            \begin{minted}[xleftmargin=-2cm, frame=leftline, fontsize=\small]{python}
                x = 7
                resultado = x + 3
            \end{minted}
        \end{itemize}
    \end{frame}

    \begin{frame}[fragile,t]{Variables}
        Calcular el área de un rectángulo
        \begin{minted}[xleftmargin=-2cm, frame=leftline, fontsize=\small]{python}
            base = 8
            altura = 5
            area = base * altura
        \end{minted}
        Calcular el promedio de tres notas
        \begin{minted}[xleftmargin=-2cm, frame=leftline, fontsize=\small]{python}
            nota1 = 4
            nota2 = 5
            nota3 = 5.2
            promedio = (nota1 + nota2 * nota3) / 3
        \end{minted}
        Expresión algebraica
        $2x + 5$
        \begin{minted}[xleftmargin=-2cm, frame=leftline, fontsize=\small]{python}
            x = 4
            resultado = 2 * x + 5
        \end{minted}
    \end{frame}

    \begin{frame}{Reglas prácticas para variables}
        \begin{itemize}
            \item Conviene usar nombres que expliquen su contenido
            \item Deben empezar con una letra ($A-Z$, $a-z$) o guion bajo (\_)
            \item Puede incluir letras, números y guiones bajos
            \item No se permiten espacios ni caracteres especiales (!, @, \#, etc)
            \item Son sensibles a mayúsculas: $edad$, $Edad$ y $EDAD$ son variables distintas
            \item Se recomienda usar snake\_case (nombre\_usuario, total\_precio)
            \item No se pueden usar palabras reservadas propias del lenguaje como \mintinline{python}{print, if}, etc
        \end{itemize}
    \end{frame}

    \begin{frame}[fragile,t]{Python - Variables}
        \begin{minted}[xleftmargin=-2cm, frame=leftline, fontsize=\small]{python}
            a = 20
            b = 5
            c = a * b
        \end{minted}
        
        El nombre de la variable también comunica
        
        \begin{minted}[xleftmargin=-2cm, frame=leftline, fontsize=\small]{python}
            precio = 20
            cantidad = 5
            total = precio * cantidad
        \end{minted}
    \end{frame}

    \begin{frame}{Identificar variables - miniactividad}
        En cada situación, identificar qué datos se deben guardar en variables
        \begin{itemize}
            \item Calcular el sueldo total de un trabajador
            \item Calcular el promedio de un estudiante
            \item Determinar si una persona puede entrar a una actividad según su edad
            \item Calcular el total a pagar por varios productos
        \end{itemize}
    \end{frame}

    \begin{frame}{Operadores}
        \begin{itemize}
            \item Permiten realizar acciones sobre los valores almacenados en variables
            \item Sirven para sumar, restar, multiplicar, concatenar texto
            \item Transforman datos
            \item Permiten obtener nuevos resultados
        \end{itemize}
    \end{frame}

    \begin{frame}[fragile,t]{Operadores aritméticos}
        \centering
        \begin{tabular}{m{0.2\textwidth} m{0.3\textwidth} m{0.4\textwidth}}
            $+$ & suma & \mintinline{python}{3 + 2 = 5} \\
            $-$ & resta & \mintinline{python}{8 - 3 = 5} \\
            $*$ & multiplicación & \mintinline{python}{4 * 2 = 8} \\
            $/$ & división & \mintinline{python}{10 / 2 = 5} \\
            $//$ & división entera & \mintinline{python}{10 // 3 = 3} \\
            $\%$ & residuo & \mintinline{python}{10 % 3 = 1} \\
            $**$ & potencia & \mintinline{python}{2 ** 3 = 8}
        \end{tabular}
    \end{frame}

    \begin{frame}[fragile,t]{Operadores}
        \begin{minted}[xleftmargin=-2cm, frame=leftline, fontsize=\small]{python}
            precio = 15
            cantidad = 3
            total = precio * cantidad
        \end{minted}
        
        \begin{minted}[xleftmargin=-2cm, frame=leftline, fontsize=\small]{python}
            saldo = 100
            gasto = 35
            nuevo_saldo = saldo - gasto
        \end{minted}
        
        \begin{minted}[xleftmargin=-2cm, frame=leftline, fontsize=\small]{python}
            galletas = 10
            personas = 3
            sobran = galletas % personas
        \end{minted}
    \end{frame}

    \begin{frame}[fragile,t]{Operadores}
        El mismo símbolo puede cambiar según el dato
        \begin{minted}[xleftmargin=-2cm, frame=leftline, fontsize=\small]{python}
            nombre = "Ana"
            apellido = "Lopez"
            completo = nombre + " " + apellido
        \end{minted}
        Aquí $+$ no suma números sino \textbf{une textos}.
    \end{frame}

    \begin{frame}{Operadores - mini actividad}
        Qué operador usarías?
        \begin{itemize}
            \item Calcular el total de 4 fusiles de 2 millones de pesos cada uno
            \item Saber cuánto dinero queda después de pagar una deuda
            \item Repartir 25 granadas entre 4 soldados y saber cuántos sobran
            \item Calcular el cuadrado de un número
        \end{itemize}
    \end{frame}

    \begin{frame}[fragile,t]{Comparadores}
        \begin{itemize}
            \item Comparar es igual a hacer una pregunta con respuesta verdadero o falso
            \item Los comparadores revisan si una condición se cumple o no
        \end{itemize}
        \centering
        \begin{minted}{python}
            edad >= 18
            nota < 4
            saldo == 0
        \end{minted}
    \end{frame}

    \begin{frame}{Comparadores}
        \begin{tabular}{m{0.3\textwidth} m{0.5\textwidth}}
            $==$ & igual que \\
            $!=$ & diferente de \\
            $>$ & mayor que \\
            $<$ & menor que \\
            $>=$ & mayor o igual que \\
            $<=$ & menor o igual que \\
        \end{tabular}
    \end{frame}

    \begin{frame}[fragile,t]{Comparadores}
        \begin{minted}[xleftmargin=-2cm, frame=leftline, fontsize=\small]{python}
            5 > 3      # True
            8 == 8     # True
            4 != 7     # True
            6 < 2      # False
        \end{minted}
    \end{frame}

    \begin{frame}[fragile,t]{Diferencia entre $=$ y $==$}
        \begin{itemize}
            \item $=$ asigna un valor a una variable
            \item $==$ compara si dos valores son iguales
        \end{itemize}
        \begin{minted}[xleftmargin=-2cm, frame=leftline, fontsize=\small]{python}
            edad = 18
            edad == 18
        \end{minted}
    \end{frame}

    \begin{frame}{Comparadores}
        \begin{itemize}
            \item $edad >= 18$ -> puede votar
            \item $nota >= 4$ -> aprueba
            \item $saldo >= precio$ -> puede comprar
            \item $temperatura > 38$ -> alerta de fiebre
            \item $intentos$ != 3 -> aún puede intentar
        \end{itemize}
    \end{frame}

    \begin{frame}[fragile,t]{Comparadores, variables y operadores}
        \begin{minted}[xleftmargin=-2cm, frame=leftline, fontsize=\small]{python}
            nota1 = 4
            nota2 = 4.5
            promedio = (nota1 + nota2) / 2
            promedio >= 4
        \end{minted}
        \begin{itemize}
            \item Las variables guardan las notas
            \item Los operadores calculan el promedio
            \item El comparador verifica si aprueba
        \end{itemize}
    \end{frame}

    \begin{frame}{Resumen}
        \begin{itemize}
            \item Las variables guardan datos
            \item Los operadores procesan esos datos
            \item Los comparadores revisan condiciones
        \end{itemize}
    \end{frame}

    
    % \end{darkframes}

    % \begin{frame}{Conclusión}
    %     \begin{itemize}
    %         \item El algoritmo es una secuencia de pasos ordenados y finitos que ayudan a resolver un problema
    %         \item El pseudocśodigo 
    %     \end{itemize}
    % \end{frame}

    % \begin{frame}[label=bibliography]{Bibliografía}
    %   \framesubtitle{\TeX, \LaTeX, and Beamer}
    %   \begin{thebibliography}{9}
    %     \bibitem{knuth84}
    %         Donald~E.~Knuth.
    %         \emph{The \TeX book}.
    %         Addison-Wesley, 1984.
    %     \bibitem{lamport94}
    %         Leslie~Lamport.
    %         \emph{\LaTeX : A Document Preparation System}.
    %         Addison-Wesley, 1986.
    %     \bibitem{MG94}
    %         M.~Goossens, F.~Mittelbach, and A.~Samarin.
    %         \emph{The \LaTeX\ Companion}.
    %         Addison-Wesley, 1994.
    %     \bibitem{tantau04}
    %         Till~Tantau.
    %         \emph{User's Guide to the Beamer Class Version 3.01}.
    %         Available at \url{http://latex-beamer.sourceforge.net}.
    %     \bibitem{MS05}
    %         A.~Mertz and W.~Slough.
    %         Edited by B.~Beeton and K.~Berry.
    %         \emph{Beamer by example} In TUGboat,
    %           Vol. 26, No. 1., pp. 68-73.
    %   \end{thebibliography}
    % \end{frame}

    % \subsection{Bibliografía}
    % \againframe{citations}
    % \againframe{bibliography}
\end{document}
